% ------------------------------------------------------------------- %
%  Pendahuluan / Introduction
%  AUTHOR: @NoHaitch IF'22
%  DATE: 2026-05-16
% ------------------------------------------------------------------- %

% -------------------- I --------------------
\chapter{Pendahuluan}
\label{1-Pendahuluan}
% -------------------- I --------------------

% -------------------- I.1 --------------------
\section{Latar Belakang}
\label{1-1-Latar-Belakang}
% -------------------- I.1 --------------------

% Note: Recommended Flow biar mudah dibaca dan dimengerti: 
%           Umum -> Spesifik -> Masalah -> Gap -> Tujuan Penelitian.

\todo{
    Bab Pendahuluan secara umum yang dijadikan landasan kerja dan arah kerja penulis tugas akhir, berfungsi mengantar pembaca untuk membaca laporan tugas akhir secara keseluruhan.
}


% -------------------- I.2 --------------------
\section{Rumusan Masalah}
\label{1-2-Rumusan-Masalah}
% -------------------- I.2 --------------------

% Note: DASAR dari seluruh penelitian. 
%       Pastikan Anda dapat menjawab pertanyaan tentang isi penelitian:
%       "Kenapa hal X ada? Apa koneksi atau relasi dengan Rumusan Masalah/Tujuan?"

\todo{
    Rumusan Masalah berisi masalah utama yang dibahas dalam tugas akhir. Rumusan masalah yang baik memiliki struktur sebagai berikut:

    \begin{enumerate}
        \item Penjelasan ringkas tentang kondisi/situasi yang ada sekarang terkait dengan topik utama yang dibahas Tugas Akhir.
        \item Pokok persoalan dari kondisi/situasi yang ada, dapat dilihat dari kelemahan atau kekurangannya. Bagian ini merupakan inti dari rumusan masalah.
        \item Elaborasi lebih lanjut yang menekankan pentingnya untuk menyelesaikan pokok persoalan tersebut.
        \item Usulan singkat terkait dengan solusi yang ditawarkan untuk menyelesaikan persoalan.
    \end{enumerate}

    Penting untuk diperhatikan bahwa persoalan yang dideskripsikan pada subbab ini akan dipertanggungjawabkan di bab Evaluasi apakah terselesaikan atau tidak.
}


% -------------------- I.3 --------------------
\section{Tujuan}
\label{1-3-Tujuan}
% -------------------- I.3 --------------------

% Note: Menjawab Rumusan Masalah, namun disarankan dengan detail atau teknik atau tambahan. 
%       Jadi bukan sekedar kopas Rumusan Masalah dan jawab.

\todo{
    Tuliskan tujuan utama dan/atau tujuan detil yang akan dicapai dalam pelaksanaan tugas akhir. Fokuskan pada hasil akhir yang ingin diperoleh setelah tugas akhir diselesaikan, terkait dengan penyelesaian persoalan pada rumusan masalah. Penting untuk diperhatikan bahwa tujuan yang dideskripsikan pada subbab ini akan dipertanggungjawabkan di akhir pelaksanaan tugas akhir apakah tercapai atau tidak.
}


% -------------------- I.4 --------------------
\section{Batasan Masalah}   
\label{1-4-Batasan-Masalah}
% -------------------- I.4 --------------------

% Note: Isi sambil berjalan saja. Pasti akan berubah-ubah.
%       Bedakan "batasan masalah" dengan "batasan implementasi"
%       e.g. Batasaan implementasi: menggunakan bahasa python.
%       e.g. Batasan masalah: penelitian hanya fokus pada modul/mekanisme X bukan pada keseluruhan sistem.

\todo{
    Tuliskan batasan masalah yang akan dibahas dalam tugas akhir. Batasan masalah ini berfungsi untuk memperjelas ruang lingkup pembahasan tugas akhir, sehingga pembaca dapat memahami dengan jelas apa yang akan dibahas dan apa yang tidak akan dibahas dalam tugas akhir. Batasan masalah ini juga berfungsi untuk menghindari pembahasan yang terlalu luas atau tidak relevan dengan topik utama tugas akhir.
}


% -------------------- I.5 --------------------
\section{Metodologi}
\label{sec:metodologi}
% -------------------- I.5 --------------------

% Note: Pilih Metodologi sebagai dasar --> Buat tahapan spesifik untuk penelitian Anda.
%       Contoh: SWE = Waterfall, Recursive, etc
%               Research = Design Science Research, etc

\todo{
    Tuliskan semua tahapan yang akan dilalui selama pelaksanaan tugas akhir. Tahapan ini spesifik untuk menyelesaikan persoalan tugas akhir. Tahapan studi literatur tidak perlu dituliskan karena ini adalah pekerjaan yang harus Anda lakukan selama proses pelaksanaan tugas akhir.
}


% -------------------- I.6 --------------------
\section{Sistematika Pembahasan}
\label{1-6-Sistematika-Pembahasan}
% -------------------- I.6 --------------------

% Note: Sebutkan semua bab pada penelitian dan jelasakan apa isi/tujuan dari bab.

\todo{
    Subbab ini berisi penjelasan ringkas isi per bab. Penjelasan ditulis satu paragraf per bab buku.
}

